\documentclass[../../main.tex]{subfiles}% 注意这里的文件路径不能用 ./main.tex ,否则用latexmk编译子文件会报错
\graphicspath{{\subfix{./image/}}{\subfix{../../image/}}} % 指定图片引用路径,后续可以直接使用图片文件名
% 如果图片存放在上级目录,则使用 ../../image/ 作为备用路径

% 例如:
% \begin{figure}[H]
% \centering
% \includegraphics[width=0.3\textwidth]{图.png}
% \caption{}
% \label{figure:图}
% \end{figure}
% 注意:上述\label{}一定要放在\caption{}之后,否则引用图片序号会只会显示??.

\begin{document}

\section{渐近性态}\label{section:ch08-渐近性态}

现考察\cref{theorem:ch08-存在性} 给出的 (P) 的整体 Young 测度值解 $(u,\nu)$ 当 $t\to\infty$ 时的渐近性态。令
\begin{align}
u_m(x,t)\triangleq u(x,t+t_m),\qquad
\nu_{m\,(x,t)}\triangleq \nu_{(x,t+t_m)} \label{align:ch08-8.68}
\end{align}
对几乎所有 $(x,t)\in Q_T$，其中 $\{t_m\}$ 为任意发散序列，并记 $X_p\equiv W^{1,p}_0(\Omega)$（$N=1$ 时 $p\in(1,\infty)$，$N\geqslant2$ 时 $p\in(1,2]$）。

\begin{definition}[$\omega$--极限点]\label{definition:ch08-omega极限点}
设 $\tilde u\in L^\infty(0,\infty;X_p)\cap C([0,\infty);L^p(\Omega))$（$N=1$ 时 $p=2$，$N\geqslant2$ 时 $p\in(1,2]$），$\tilde\nu\in\mathcal{Y}^+(Q_\infty;\mathbb{R}^N)$。偶 $(\tilde u,\tilde\nu)$ 称为问题 (P) 的整体 Young 测度值解 $(u,\nu)$ 的 $\omega$--极限点，若：
\begin{enumerate}[(i)]
\item\label{item:ch08-omega-i} 存在发散序列 $\{t_m\}\subseteq(0,\infty)$，使
\begin{enumerate}[a)]
\item\label{item:ch08-omega-i-a} $u_m\rightharpoonup\tilde u$ 在 $L^r(0,T;X_p)$（对所有 $T\in(0,\infty)$ 与 $r\in[1,\infty)$），
\item\label{item:ch08-omega-i-b} $\nu_m\stackrel{*}{\rightharpoonup}\tilde\nu$ 在 $\mathcal{Y}^+(Q_T;\mathbb{R}^N)$（对所有 $T\in(0,\infty)$）；
\end{enumerate}
\item\label{item:ch08-omega-ii} $(\tilde u,\tilde\nu)$ 是问题 (P) 以函数 $\tilde u(\cdot,0)$ 为初值的整体 Young 测度值解。
\end{enumerate}
$(u,\nu)$ 的 $\omega$--极限点的集合称为 $(u,\nu)$ 的 $\omega$--极限集。
\end{definition}

我们证明如下结果。

\begin{theorem}[$\omega$--极限集的存在性与支撑]\label{theorem:ch08-omega极限非空}
设或 $N=1$ 且\eqref{align:ch08-8.1a}--\eqref{align:ch08-8.1d} 成立，或 $N\geqslant2$ 且\eqref{align:ch08-8.2a}--\eqref{align:ch08-8.2f} 成立。设 $(u,\nu)\in L^\infty(0,\infty;X_p)\times\mathcal{Y}^+(Q_\infty;\mathbb{R}^N)$ 为问题 (P) 的整体 Young 测度值解（其存在性由\cref{theorem:ch08-存在性} 保证）。则：
\begin{enumerate}[(i)]
\item\label{item:ch08-omega-thm-i} $(u,\nu)$ 的 $\omega$--极限集非空；
\item\label{item:ch08-omega-thm-ii} 若 $u_0\in H^1_0(\Omega)$，且 $N\geqslant2$ 时\eqref{align:ch08-8.2a}--\eqref{align:ch08-8.2f} 取 $p=2$ 成立，则对每个 $\omega$--极限点 $(\tilde u,\tilde\nu)$，
\begin{align}
\operatorname{supp}\tilde\nu_{(x,t)}\subseteq S\triangleq \{y\in\mathbb{R}^N\mid\phi(y)\cdot y=0\} \label{align:ch08-8.69}
\end{align}
对几乎所有 $(x,t)\in Q_\infty$。
\end{enumerate}
\end{theorem}

\begin{theorem}[$\omega$--极限为平衡解]\label{theorem:ch08-平衡解}
设\cref{theorem:ch08-omega极限非空} 的假设满足，$u_0\in H^1_0(\Omega)$，且 $N\geqslant2$ 时\eqref{align:ch08-8.2a}--\eqref{align:ch08-8.2f} 取 $p=2$ 成立。若 $S=\{y\in\mathbb{R}^N\mid\phi(y)=0\}$，则 $(u,\nu)$ 的每个 $\omega$--极限点 $(\tilde u,\tilde\nu)$ 都是问题 (P) 的 Young 测度平衡解。
\end{theorem}

为证\cref{theorem:ch08-omega极限非空}，需要一些预备结果。

\begin{lemma}[平移序列 $u_m$ 的收敛]\label{lemma:ch08-u_m收敛}
设\cref{theorem:ch08-omega极限非空} 的假设满足。则存在\eqref{align:ch08-8.68} 中序列 $\{u_m\}$ 的子列（为简便不重新标号）与 $\tilde u\in L^\infty(0,\infty;X_p)\cap C([0,\infty);L^p(\Omega))\cap W^{1,p}(Q_T)$（$p\in(1,2]$），$\partial_t\tilde u\in L^2(Q_\infty)$，使对所有 $T\in(0,\infty)$，
\begin{align}
&u_m\to\tilde u\ \text{在 }C([0,T);L^p(\Omega))\ (p\in(1,2]); \label{align:ch08-8.70a}\\
&u_m\rightharpoonup\tilde u\ \text{在 }H^1(Q_T)\ (N=1),\quad
u_m\rightharpoonup\tilde u\ \text{在 }W^{1,p}(Q_T)\ (N\geqslant2); \label{align:ch08-8.70b}\\
&u_m\rightharpoonup\tilde u\ \text{在 }L^r(0,T;X_p)\ (r\in[1,\infty)); \label{align:ch08-8.70c}\\
&u_m\stackrel{*}{\rightharpoonup}\tilde u\ \text{在 }L^\infty(0,T;X_p); \label{align:ch08-8.70d}\\
&\nabla u_m\stackrel{*}{\rightharpoonup}\nabla\tilde u\ \text{在 }L^\infty(0,T;L^p(\Omega))\ (p\in(1,2]). \label{align:ch08-8.70e}
\end{align}
进而
\begin{align}
\partial_t u_m\rightharpoonup\partial_t\tilde u\ \text{在 }L^2(Q_\infty). \label{align:ch08-8.70f}
\end{align}
\end{lemma}

\begin{proof}
回忆 $u\in L^\infty(0,\infty;X_p)$、$\partial_t u\in L^2(Q_\infty)$；且对所有 $T\in(0,\infty)$ 有 $u\in W^{1,p}(Q_T)\cap C([0,\infty);L^p(\Omega))$（$p\in(1,2]$；见\cref{proposition:ch08-收敛N2} 与\cref{proposition:ch08-收敛N1} 及\cref{theorem:ch08-存在性}）。显然，对所有 $n\in\mathbb{N}$，
\begin{align}
\|u_m\|_{L^\infty(0,\infty;X_p)}+\|\partial_t u_m\|_{L^2(Q_\infty)}
\leqslant\|u\|_{L^\infty(0,\infty;X_p)}+\|\partial_t u\|_{L^2(Q_\infty)}, \label{align:ch08-8.71}\\
\|u_m\|_{H^1(Q_T)}\leqslant\{T\|u\|^2_{L^\infty(\mathbb{R}^+;X_p)}+\|u_t\|^2_{L^2(Q_\infty)}\}^{1/2}\quad(N=1), \label{align:ch08-8.72}\\
\|u_m\|_{W^{1,p}(Q_T)}\leqslant\{T\|u\|^p_{L^\infty(\mathbb{R}^+;X_p)}+(\lambda(\Omega)T)^{1-p/2}\|u_t\|^p_{L^2(Q_\infty)}\}^{1/p}\quad(N\geqslant2,\ p\in(1,2]) \label{align:ch08-8.73}
\end{align}
对所有 $T\in(0,\infty)$。如\cref{proposition:ch08-收敛N2} 证明那样，\eqref{align:ch08-8.70a}--\eqref{align:ch08-8.70b} 中的收敛由不等式\eqref{align:ch08-8.72}--\eqref{align:ch08-8.73} 得到，而\eqref{align:ch08-8.70c}--\eqref{align:ch08-8.70f} 中的收敛由不等式\eqref{align:ch08-8.71} 得到。结论即得。
\end{proof}

\begin{lemma}[平移 Young 测度列的紧性]\label{lemma:ch08-紧性}
设\cref{theorem:ch08-omega极限非空} 的假设满足。则\eqref{align:ch08-8.68} 定义的平移 Young 测度序列 $\{\nu_m\}\subseteq\mathcal{Y}^+(Q_\infty;\mathbb{R}^N)$ 在 $\mathcal{Y}^+(Q_T;\mathbb{R}^N)$ 中对所有 $T\in(0,\infty)$ 紧。
\end{lemma}

\begin{proof}
固定 $T\in(0,\infty)$。由估计\eqref{align:ch08-8.21} 与\eqref{align:ch08-8.29}，序列 $\{|\nabla u_{\epsilon_k}|\}$ 在 $L^p(Q_\infty)$ 中有界（$N=1$ 时 $p\in(1,\infty)$，$N\geqslant2$ 时 $p\in(1,2]$），因而由\cref{lemma:ch02-有界一致可积} 在每个 $Q_\tau=\Omega\times(0,\tau)$（$\tau>0$）中对每个 $q\in(1,p)$ 为 $q$--一致可积。于是由\cref{proposition:ch05-弱收敛}，对每个 $m\in\mathbb{N}$，
\begin{align}
C\geqslant\lim_{k\to\infty}\iint_{Q_{m,T}}|\nabla u_{\epsilon_k}|^q\,\mathrm{d}x\,\mathrm{d}t
&=\iint_{Q_{m,T}}\mathrm{d}x\,\mathrm{d}t\int_{\mathbb{R}^N}|y|^q\,\mathrm{d}\nu_{(x,t)}(y)\nonumber\\
&=\iint_{Q_T}\mathrm{d}x\,\mathrm{d}t\int_{\mathbb{R}^N}|y|^q\,\mathrm{d}\nu_{m\,(x,t)}(y) \label{align:ch08-8.74}
\end{align}
其中常数 $C>0$，$Q_{m,T}\triangleq \Omega\times(t_m,t_m+T)$。令 $K_j\triangleq \{y\in\mathbb{R}^N\mid|y|\leqslant j\}$，它紧；在余集 $K_j^c$ 上有 $|y|>j$（$j\in\mathbb{N}$）。由 Chebyshev 不等式与\eqref{align:ch08-8.74} 有
\begin{align*}
j^q\nu_m(Q_T\times K_j^c)\leqslant\iint_{Q_T}\mathrm{d}x\,\mathrm{d}t\int_{\mathbb{R}^N}|y|^q\,\mathrm{d}\nu_{m\,(x,t)}(y)\leqslant C,
\end{align*}
因而
\begin{align}
\nu_m(Q_T\times K_j^c)\leqslant\frac{C}{j^q}\quad\text{对所有 }j,m\in\mathbb{N}. \label{align:ch08-8.75}
\end{align}
任取 $\epsilon>0$，选 $j_0>(C/\epsilon)^{1/q}$。由\eqref{align:ch08-8.75} 有 $\nu_m(Q_T\times K_{j_0}^c)<\epsilon$ 对所有 $m\in\mathbb{N}$，结论即得。
\end{proof}

\begin{lemma}[平移 Young 测度列的收敛]\label{lemma:ch08-ν_m收敛}
设\cref{theorem:ch08-omega极限非空} 的假设满足。则存在\eqref{align:ch08-8.68} 中序列 $\{\nu_m\}$ 的子列（为简便不重新标号）与 $\tilde\nu\in\mathcal{Y}^+(Q_\infty;\mathbb{R}^N)$，使 $\nu_m\stackrel{*}{\rightharpoonup}\tilde\nu$ 在 $\mathcal{Y}^+(Q_T;\mathbb{R}^N)$ 对所有 $T\in(0,\infty)$。
\end{lemma}

\begin{proof}
由\cref{lemma:ch08-紧性} 与\cref{theorem:ch05-弱紧Young} 经标准对角论证即得。
\end{proof}

\begin{lemma}[平移测度的积分收敛]\label{lemma:ch08-积分收敛}
设\cref{theorem:ch08-omega极限非空} 的假设满足。设 $\{\nu_m\}\subseteq\mathcal{Y}^+(Q_\infty;\mathbb{R}^N)$ 为\eqref{align:ch08-8.68} 定义的序列，$\rho\in C(\mathbb{R}^N)$ 满足 $|\rho(y)|\leqslant M(1+|y|^q)$（对某个 $M>0$ 与 $q\in(0,p)$，其中 $N=1$ 时 $p\in(1,\infty)$，$N\geqslant2$ 时 $p\in(1,2]$）。则对每个 $T\in(0,\infty)$：
\begin{enumerate}[(i)]
\item\label{item:ch08-intconv-i} 映射 $(x,t)\mapsto\int_{\mathbb{R}^N}\rho(y)\,\mathrm{d}\nu_{m\,(x,t)}(y)$ 属于 $L^1(Q_T)$；
\item\label{item:ch08-intconv-ii} 对所有 $\zeta\in L^\infty(Q_T)$，
\begin{align}
\lim_{j\to\infty}\iint_{Q_T}\zeta(x,t)\,\mathrm{d}x\,\mathrm{d}t\int_{\mathbb{R}^N}f_j(|y|)\rho(y)\,\mathrm{d}\nu_{m\,(x,t)}(y)
=\iint_{Q_T}\zeta(x,t)\,\mathrm{d}x\,\mathrm{d}t\int_{\mathbb{R}^N}\rho(y)\,\mathrm{d}\nu_{m\,(x,t)}(y) \label{align:ch08-8.76}
\end{align}
其中 $f_j\in C_c([0,\infty))$ 如\eqref{align:ch08-8.56}，且该收敛关于 $m\in\mathbb{N}$ 一致。
\end{enumerate}
\end{lemma}

\begin{proof}
(i) 由不等式\eqref{align:ch08-8.21} 与\eqref{align:ch08-8.29}，经\cref{lemma:ch02-有界一致可积}，对每个 $q\in(0,p)$，序列 $\{|\nabla u_{\epsilon_k}|^q\}\subseteq L^1(Q_{m,T})$ 有界且一致可积。由 $\rho$ 的增长假设与\cref{proposition:ch05-弱收敛} 有
\begin{align*}
\iint_{Q_T}\mathrm{d}x\,\mathrm{d}t\left|\int_{\mathbb{R}^N}\rho(y)\,\mathrm{d}\nu_{m\,(x,t)}(y)\right|
&\leqslant M\iint_{Q_T}\mathrm{d}x\,\mathrm{d}t\int_{\mathbb{R}^N}(1+|y|^q)\,\mathrm{d}\nu_{m\,(x,t)}(y)\\
&=M\iint_{Q_{m,T}}\mathrm{d}x\,\mathrm{d}t\int_{\mathbb{R}^N}(1+|y|^q)\,\mathrm{d}\nu_{(x,t)}(y)\\
&=M\lim_{k\to\infty}\iint_{Q_{m,T}}(1+|\nabla u_{\epsilon_k}|^q)\,\mathrm{d}x\,\mathrm{d}t<\infty
\end{align*}
（可能抽 $\{u_{\epsilon_k}\}$ 的子列，不重新标号）。结论即得。

(ii) 令 $g_j\triangleq 1-f_j$；则 $0\leqslant g_j\leqslant1$ 在 $[0,\infty)$ 上且 $g_j=0$ 在 $[0,j]$ 上（$j\in\mathbb{N}$）。由于序列 $\{|\nabla u_{\epsilon_k}|^q\}\subseteq L^1(Q_{m,T})$ 有界且一致可积（见 (i)），由 $\rho$ 的增长假设，每个序列 $\{g_j(|\nabla u_{\epsilon_k}|)|\rho(\nabla u_{\epsilon_k})|\}$（$j\in\mathbb{N}$）亦然。于是由\cref{proposition:ch05-弱收敛}，存在 $\{g_j(|\nabla u_{\epsilon_k}|)|\rho(\nabla u_{\epsilon_k})|\}$ 的子列（不重新标号，可依赖于 $j,m$），使对所有 $\zeta\in L^\infty(Q_T)$，
\begin{align}
\iint_{Q_{m,T}}\zeta(x,t)\,\mathrm{d}x\,\mathrm{d}t\int_{\mathbb{R}^N}g_j(|y|)|\rho(y)|\,\mathrm{d}\nu_{(x,t)}(y)
=\lim_{k\to\infty}\iint_{Q_{m,T}}[g_j(|\nabla u_{\epsilon_k|})|\rho(\nabla u_{\epsilon_k})|](x,t)\zeta(x,t)\,\mathrm{d}x\,\mathrm{d}t \label{align:ch08-8.77}
\end{align}
（$j,m\in\mathbb{N}$）。于是对所有 $\zeta\in L^\infty(Q_T)$ 与 $j,m\in\mathbb{N}$，
\begin{align}
\left|\iint_{Q_{m,T}}\zeta(x,t)\,\mathrm{d}x\,\mathrm{d}t\int_{\mathbb{R}^N}g_j(|y|)\rho(y)\,\mathrm{d}\nu_{(x,t)}(y)\right|
\leqslant\|\zeta\|_{L^\infty(Q_T)}\sup_{k\in\mathbb{N}}\iint_{Q_{m,T}}[g_j(|\nabla u_{\epsilon_k|})|\rho(\nabla u_{\epsilon_k})|](x,t)\,\mathrm{d}x\,\mathrm{d}t. \label{align:ch08-8.78}
\end{align}
由 $\rho$ 的增长条件存在 $M_0>0$，使对所有 $|y|\geqslant1$ 有 $|\rho(y)|\leqslant M_0|y|^q$。令 $E_{j,k,m}\triangleq \{(x,t)\in Q_{m,T}\mid|\nabla u_{\epsilon_k}|(x,t)\geqslant j\}$；于是 $g_j(|\nabla u_{\epsilon_k}|)=0$ 在 $(E_{j,k,m})^c$ 上，$g_j(|\nabla u_{\epsilon_k}|)\leqslant1$ 在 $E_{j,k,m}$ 上。对任意固定的 $j,k,m\in\mathbb{N}$，有
\begin{align}
&\iint_{Q_{m,T}}[g_j(|\nabla u_{\epsilon_k|})|\rho(\nabla u_{\epsilon_k})|](x,t)\,\mathrm{d}x\,\mathrm{d}t\nonumber\\
&\quad=\iint_{E_{j,k,m}}[g_j(|\nabla u_{\epsilon_k|})|\rho(\nabla u_{\epsilon_k})|](x,t)\,\mathrm{d}x\,\mathrm{d}t\nonumber\\
&\quad\leqslant M_0\iint_{E_{j,k,m}}|\nabla u_{\epsilon_k}|^q(x,t)\,\mathrm{d}x\,\mathrm{d}t\nonumber\\
&\quad\leqslant M_0\left(\iint_{E_{j,k,m}}|\nabla u_{\epsilon_k}|^p(x,t)\,\mathrm{d}x\,\mathrm{d}t\right)^{q/p}
[\lambda_{N+1}(E_{j,k,m})]^{1-q/p}. \label{align:ch08-8.79}
\end{align}
另一方面，由 Chebyshev 不等式有
\begin{align*}
\lambda_{N+1}(E_{j,k,m})\leqslant\left(\iint_{E_{j,k,m}}|\nabla u_{\epsilon_k}|^p(x,t)\,\mathrm{d}x\,\mathrm{d}t\right)j^{-p}.
\end{align*}
由\eqref{align:ch08-8.79} 有
\begin{align}
&\iint_{Q_{m,T}}[g_j(|\nabla u_{\epsilon_k|})|\rho(\nabla u_{\epsilon_k})|](x,t)\,\mathrm{d}x\,\mathrm{d}t\nonumber\\
&\quad\leqslant M_0\left(\iint_{E_{j,k,m}}|\nabla u_{\epsilon_k}|^p(x,t)\,\mathrm{d}x\,\mathrm{d}t\right)j^{-(p-q)}
\leqslant KM_0T\,j^{-(p-q)} \label{align:ch08-8.80}
\end{align}
其中常数 $K>0$ 仅依赖于 $u_0$ 与 $\lambda_N(\Omega)$（见不等式\eqref{align:ch08-8.21} 与\eqref{align:ch08-8.29} 的证明）。由\eqref{align:ch08-8.78} 与\eqref{align:ch08-8.80} 得
\begin{align*}
\left|
\begin{aligned}
&\iint_{Q_T}\zeta(x,t)\,\mathrm{d}x\,\mathrm{d}t\int_{\mathbb{R}^N}\rho(y)\,\mathrm{d}\nu_{m\,(x,t)}(y)
-\iint_{Q_T}\zeta(x,t)\,\mathrm{d}x\,\mathrm{d}t\int_{\mathbb{R}^N}f_j(|y|)\rho(y)\,\mathrm{d}\nu_{m\,(x,t)}(y)\\
&=\left|\iint_{Q_{m,T}}\zeta(x,t)\,\mathrm{d}x\,\mathrm{d}t\int_{\mathbb{R}^N}g_j(|y|)\rho(y)\,\mathrm{d}\nu_{(x,t)}(y)\right|
\leqslant KM_0T\,j^{-(p-q)}\|\zeta\|_\infty,
\end{aligned}
\right.
\end{align*}
故断言 (ii) 得证。证毕。
\end{proof}

现可证\cref{theorem:ch08-omega极限非空}。

\begin{proof}[\cref{theorem:ch08-omega极限非空} 的证明]
(i) 设 $(\tilde u,\tilde\nu)$ 为\cref{lemma:ch08-u_m收敛} 与\cref{lemma:ch08-ν_m收敛} 中相应分量的偶。由这些引理，$(\tilde u,\tilde\nu)$ 满足\cref{definition:ch08-omega极限点} 的要求 (i)。于是若证明了 $(\tilde u,\tilde\nu)$ 还满足\cref{definition:ch08-Young测度值解} 的等式\eqref{align:ch08-8.4} 与\eqref{align:ch08-8.5}（Cauchy 数据为 $\tilde u(\cdot,0)$），断言 (i) 即得。

先证 $(\tilde u,\tilde\nu)$ 满足等式\eqref{align:ch08-8.4}，即
\begin{align}
\nabla\tilde u(x,t)=\int_{\mathbb{R}^N}y\,\mathrm{d}\tilde\nu_{(x,t)}(y)
\quad\text{对几乎所有 }(x,t)\in Q_T, \label{align:ch08-8.81}
\end{align}
其中 $\tilde\nu_{(x,t)}\in\mathcal{P}(\mathbb{R}^N)$ 表示 $\tilde\nu$ 的离散分解。回忆定义\eqref{align:ch08-8.68}，由等式\eqref{align:ch08-8.4} 得
\begin{align}
\nabla u_m(x,t)=\int_{\mathbb{R}^N}y\,\mathrm{d}\nu_{m\,(x,t)}(y) \label{align:ch08-8.82}
\end{align}
对几乎所有 $(x,t)\in Q_T$ 与所有 $m\in\mathbb{N}$。由\eqref{align:ch08-8.70c} 中的收敛，对所有 $\zeta\in L^\infty(Q_T)$ 有
\begin{align}
\lim_{m\to\infty}\iint_{Q_T}\nabla u_m(x,t)\zeta(x,t)\,\mathrm{d}x\,\mathrm{d}t
=\iint_{Q_T}\nabla\tilde u(x,t)\zeta(x,t)\,\mathrm{d}x\,\mathrm{d}t. \label{align:ch08-8.83}
\end{align}
设 $\{f_j\}\subseteq C_c([0,\infty))$（$j\in\mathbb{N}$）满足\eqref{align:ch08-8.56}，$\zeta\in L^\infty(Q_T)$ 任意取定。由\cref{lemma:ch08-积分收敛}(ii) 有
\begin{align}
\lim_{j\to\infty}\iint_{Q_T}\zeta(x,t)\,\mathrm{d}x\,\mathrm{d}t\int_{\mathbb{R}^N}f_j(|y|)y\,\mathrm{d}\nu_{m\,(x,t)}(y)
=\iint_{Q_T}\zeta(x,t)\,\mathrm{d}x\,\mathrm{d}t\int_{\mathbb{R}^N}y\,\mathrm{d}\nu_{m\,(x,t)}(y) \label{align:ch08-8.84}
\end{align}
关于 $m\in\mathbb{N}$ 一致。另一方面，由于函数 $h_j(x,t,y)\triangleq f_j(|y|)y\,\zeta(x,t)\in C_b(Q_T\times\mathbb{R}^N)$，由\cref{lemma:ch08-紧性} 得
\begin{align}
\lim_{m\to\infty}\iint_{Q_T}\zeta(x,t)\,\mathrm{d}x\,\mathrm{d}t\int_{\mathbb{R}^N}f_j(|y|)y\,\mathrm{d}\nu_{m\,(x,t)}(y)
=\iint_{Q_T}\zeta(x,t)\,\mathrm{d}x\,\mathrm{d}t\int_{\mathbb{R}^N}f_j(|y|)y\,\mathrm{d}\tilde\nu_{(x,t)}(y). \label{align:ch08-8.85}
\end{align}
由\eqref{align:ch08-8.82} 与\eqref{align:ch08-8.84}--\eqref{align:ch08-8.85}，
\begin{align}
\begin{aligned}
\lim_{m\to\infty}\iint_{Q_T}\nabla u_m(x,t)\zeta(x,t)\,\mathrm{d}x\,\mathrm{d}t
&=\lim_{m\to\infty}\iint_{Q_T}\zeta(x,t)\,\mathrm{d}x\,\mathrm{d}t\int_{\mathbb{R}^N}y\,\mathrm{d}\nu_{m\,(x,t)}(y)\\
&=\lim_{m\to\infty}\lim_{j\to\infty}\iint_{Q_T}\zeta(x,t)\,\mathrm{d}x\,\mathrm{d}t\int_{\mathbb{R}^N}f_j(|y|)y\,\mathrm{d}\nu_{m\,(x,t)}(y)\\
&=\lim_{j\to\infty}\lim_{m\to\infty}\iint_{Q_T}\zeta(x,t)\,\mathrm{d}x\,\mathrm{d}t\int_{\mathbb{R}^N}f_j(|y|)y\,\mathrm{d}\nu_{m\,(x,t)}(y)\\
&=\lim_{j\to\infty}\iint_{Q_T}\zeta(x,t)\,\mathrm{d}x\,\mathrm{d}t\int_{\mathbb{R}^N}f_j(|y|)y\,\mathrm{d}\tilde\nu_{(x,t)}(y)\\
&=\iint_{Q_T}\zeta(x,t)\,\mathrm{d}x\,\mathrm{d}t\int_{\mathbb{R}^N}y\,\mathrm{d}\tilde\nu_{(x,t)}(y),
\end{aligned} \label{align:ch08-8.86}
\end{align}
其中最后一等式由控制收敛定理得到，因 $\tilde\nu_{(x,t)}(\mathbb{R}^N)=1$ 对几乎所有 $(x,t)\in Q_T$，且 $\tilde\nu$ 的重心属于 $L^1(Q_T)$。由 $\zeta$ 的任意性，从\eqref{align:ch08-8.83}--\eqref{align:ch08-8.86} 得\eqref{align:ch08-8.81}。

类似地可证 $(\tilde u,\tilde\nu)$ 以初值 $\tilde u(\cdot,0)$ 满足\eqref{align:ch08-8.5}，即
\begin{align}
&\int_0^t\int_\Omega\left\{\tilde u(x,s)\partial_s\zeta(x,s)-\sum_{i=1}^N\int_{\mathbb{R}^N}\phi_i(y)\,\mathrm{d}\tilde\nu_{(x,s)}(y)\,\partial_{x_i}\zeta(x,s)\right\}\mathrm{d}x\,\mathrm{d}s\nonumber\\
&\qquad=\int_\Omega\tilde u(x,t)\zeta(x,t)\,\mathrm{d}x-\int_\Omega\tilde u(x,0)\zeta(x,0)\,\mathrm{d}x \label{align:ch08-8.87}
\end{align}
对所有 $\zeta\in C^1([0,T);C^1_c(\Omega))$ 与 $t\in(0,T)$（$T\in(0,\infty)$）。按记号滥用，仍记 $\zeta$ 的任一延拓 $\zeta\in C^1((-\infty,T);C^1_c(\Omega))$；于是对任意序列 $\{t_m\}\subseteq(0,\infty)$ 与 $(x,t)\in Q_T$，平移 $\zeta_m(x,t)\triangleq \zeta(x,t-t_m)$ 有定义且 $\zeta_m\in C^1([0,T);C^1_c(\Omega))$。由\eqref{align:ch08-8.5} 得
\begin{align*}
\int_{t_m}^{t+t_m}\int_\Omega\left\{u(x,s)\partial_s\zeta(x,s-t_m)-\sum_{i=1}^N\int_{\mathbb{R}^N}\phi_i(y)\,\mathrm{d}\nu_{(x,s)}(y)\,\partial_{x_i}\zeta(x,s-t_m)\right\}\mathrm{d}x\,\mathrm{d}s
=\int_\Omega u(x,t+t_m)\zeta(x,t)\,\mathrm{d}x-\int_\Omega u(x,t_m)\zeta(x,0)\,\mathrm{d}x
\end{align*}
对每个 $\{t_m\}\subseteq(0,\infty)$ 与所有 $t\in(0,T)$、$T\in(0,\infty)$（$0<t_m<t+t_m<T$）。上式写为
\begin{align}
&\int_0^t\int_\Omega\left\{u_m(x,s)\partial_s\zeta(x,s)-\sum_{i=1}^N\int_{\mathbb{R}^N}\phi_i(y)\,\mathrm{d}\nu_{m\,(x,s)}(y)\,\partial_{x_i}\zeta(x,s)\right\}\mathrm{d}x\,\mathrm{d}s\nonumber\\
&\qquad=\int_\Omega u_m(x,t)\zeta(x,t)\,\mathrm{d}x-\int_\Omega u_m(x,0)\zeta(x,0)\,\mathrm{d}x. \label{align:ch08-8.88}
\end{align}
由\eqref{align:ch08-8.70a} 与\eqref{align:ch08-8.70c} 中的收敛有
\begin{align}
\lim_{m\to\infty}\int_0^t\int_\Omega u_m(x,s)\partial_s\zeta(x,s)\,\mathrm{d}x\,\mathrm{d}s
=\int_0^t\int_\Omega\tilde u(x,s)\partial_s\zeta(x,s)\,\mathrm{d}x\,\mathrm{d}s, \label{align:ch08-8.89a}\\
\lim_{m\to\infty}\int_\Omega u_m(x,t)\zeta(x,t)\,\mathrm{d}x
=\int_\Omega\tilde u(x,t)\zeta(x,t)\,\mathrm{d}x, \label{align:ch08-8.89b}\\
\lim_{m\to\infty}\int_\Omega u_m(x,0)\zeta(x,0)\,\mathrm{d}x
=\int_\Omega\tilde u(x,0)\zeta(x,0)\,\mathrm{d}x. \label{align:ch08-8.89c}
\end{align}
如前，设 $\{f_j\}\subseteq C_c([0,\infty))$ 满足\eqref{align:ch08-8.56}。由\cref{lemma:ch08-积分收敛} 有
\begin{align}
\lim_{j\to\infty}\int_0^t\int_\Omega\partial_{x_i}\zeta(x,s)\,\mathrm{d}x\,\mathrm{d}s\int_{\mathbb{R}^N}f_j(|y|)\phi_i(y)\,\mathrm{d}\nu_{m\,(x,s)}(y)
=\int_0^t\int_\Omega\partial_{x_i}\zeta(x,s)\,\mathrm{d}x\,\mathrm{d}s\int_{\mathbb{R}^N}\phi_i(y)\,\mathrm{d}\nu_{m\,(x,s)}(y) \label{align:ch08-8.90}
\end{align}
关于 $m\in\mathbb{N}$ 一致。仿\eqref{align:ch08-8.86} 的论证显然给出
\begin{align}
\lim_{m\to\infty}\int_0^t\int_\Omega\partial_{x_i}\zeta(x,s)\,\mathrm{d}x\,\mathrm{d}s\int_{\mathbb{R}^N}\phi_i(y)\,\mathrm{d}\nu_{m\,(x,s)}(y)
=\int_0^t\int_\Omega\partial_{x_i}\zeta(x,s)\,\mathrm{d}x\,\mathrm{d}s\int_{\mathbb{R}^N}\phi_i(y)\,\mathrm{d}\tilde\nu_{(x,s)}(y) \label{align:ch08-8.91}
\end{align}
（$i=1,\cdots ,N$）。由\eqref{align:ch08-8.89a} 与\eqref{align:ch08-8.91} 得等式\eqref{align:ch08-8.87}。断言 (i) 得证。

(ii) 设 $f\in C([0,\infty))$ 满足 $z\leqslant1$ 时 $f(z)=1$，$z\in(1,\infty)$ 时 $0<f(z)<1$，且 $z\to\infty$ 时 $f(z)\to0$ 足够快，使函数 $h:\mathbb{R}^N\to(0,\infty)$、$h(y)\triangleq f(|y|)\phi(y)\cdot y$ 属于 $C_0(\mathbb{R}^N)$。由于对每个 $y\in\mathbb{R}^N$ 有 $f(|y|)>0$，
\begin{align*}
S\triangleq \{y\in\mathbb{R}^N\mid\phi(y)\cdot y=0\}=\{y\in\mathbb{R}^N\mid h(y)=0\}.
\end{align*}
于是若证明了
\begin{align}
\iint_{Q_T}\mathrm{d}x\,\mathrm{d}t\int_{\mathbb{R}^N}h(y)\,\mathrm{d}\tilde\nu_{(x,t)}(y)=0 \label{align:ch08-8.92}
\end{align}
结论即得（见不等式\eqref{align:ch08-8.1d} 与\eqref{align:ch08-8.2f}）。

为此，注意 $h\in C_0(\mathbb{R}^N)\subseteq C_b(Q_T\times\mathbb{R}^N)$，由\cref{lemma:ch08-紧性} 有
\begin{align}
\iint_{Q_T}\mathrm{d}x\,\mathrm{d}t\int_{\mathbb{R}^N}h(y)\,\mathrm{d}\tilde\nu_{(x,t)}(y)
=\lim_{m\to\infty}\iint_{Q_T}\mathrm{d}x\,\mathrm{d}t\int_{\mathbb{R}^N}h(y)\,\mathrm{d}\nu_{m\,(x,t)}(y). \label{align:ch08-8.93}
\end{align}
又因 $h$ 在 $\mathbb{R}^N$ 上有界，由\cref{lemma:ch02-有界一致可积} 序列 $\{h(\nabla u_{\epsilon_k})\}\subseteq L^1(Q_T)$（$T\in(0,\infty)$）有界且一致可积。由\cref{proposition:ch05-弱收敛}，对每个 $m\in\mathbb{N}$，
\begin{align}
\iint_{Q_T}\mathrm{d}x\,\mathrm{d}t\int_{\mathbb{R}^N}h(y)\,\mathrm{d}\nu_{m\,(x,t)}(y)
=\lim_{k\to\infty}\iint_{Q_{m,T}}h(\nabla u_{\epsilon_k})(x,t)\,\mathrm{d}x\,\mathrm{d}t. \label{align:ch08-8.94}
\end{align}
由\eqref{align:ch08-8.93}--\eqref{align:ch08-8.94} 得
\begin{align}
\iint_{Q_T}\mathrm{d}x\,\mathrm{d}t\int_{\mathbb{R}^N}h(y)\,\mathrm{d}\tilde\nu_{(x,t)}(y)
=\lim_{m\to\infty}\lim_{k\to\infty}\iint_{Q_{m,T}}h(\nabla u_{\epsilon_k})(x,t)\,\mathrm{d}x\,\mathrm{d}t. \label{align:ch08-8.95}
\end{align}
现由\cref{remark:ch08-能量衰减} 注意
\begin{align*}
0&\leqslant\iint_{Q_{m,T}}h(\nabla u_{\epsilon_k})(x,t)\,\mathrm{d}x\,\mathrm{d}t
\leqslant\iint_{Q_{m,T}}[\phi(\nabla u_{\epsilon_k})\cdot\nabla u_{\epsilon_k}](x,t)\,\mathrm{d}x\,\mathrm{d}t\\
&=-\frac12\Big\{
\|u_{\epsilon_k}(\cdot,T+t_m)\|^2_{L^2(\Omega)}-\|u_{\epsilon_k}(\cdot,t_m)\|^2_{L^2(\Omega)}
+\epsilon_k\big(\|\nabla u_{\epsilon_k}(\cdot,T+t_m)\|^2_{L^2(\Omega)}-\|\nabla u_{\epsilon_k}(\cdot,t_m)\|^2_{L^2(\Omega)}\big)
\Big\}.
\end{align*}
在上述不等式中令 $k\to\infty$，用\eqref{align:ch08-8.36a} 中的收敛与 $N\geqslant2$ 时取 $p=2$ 的估计\eqref{align:ch08-8.29}，或 $N=1$ 时分别用\eqref{align:ch08-8.42a} 与估计\eqref{align:ch08-8.21}（见\cref{remark:ch08-能量衰减}），由等式\eqref{align:ch08-8.94} 得
\begin{align*}
0\leqslant\iint_{Q_T}\mathrm{d}x\,\mathrm{d}t\int_{\mathbb{R}^N}h(y)\,\mathrm{d}\nu_{m\,(x,t)}(y)
\leqslant-\frac12\big[\|u(\cdot,T+t_m)\|^2_{L^2(\Omega)}-\|u(\cdot,t_m)\|^2_{L^2(\Omega)}\big]
\end{align*}
对每个 $m\in\mathbb{N}$。由于映射 $t\mapsto\|u(\cdot,t)\|^2_{L^2(\Omega)}$ 在 $(0,\infty)$ 上非增（见\cref{remark:ch08-能量衰减}），极限 $\lim_{t\to\infty}\|u(\cdot,t)\|^2_{L^2(\Omega)}\in\mathbb{R}$ 存在。在上式中令 $m\to\infty$ 并用\eqref{align:ch08-8.93}，即得\eqref{align:ch08-8.92}。证毕。
\end{proof}

\begin{proof}[\cref{theorem:ch08-平衡解} 的证明]
由\cref{theorem:ch08-omega极限非空}(ii)，在当前假设下对几乎所有 $(x,t)\in Q_T$ 有 $\operatorname{supp}\tilde\nu_{(x,t)}\subseteq\{\phi=0\}$，因而等式\eqref{align:ch08-8.7} 满足。于是对所有 $\rho\in C^1_c(\Omega)$，由\eqref{align:ch08-8.87} 得
\begin{align*}
\int_\Omega\tilde u(x,t)\rho(x)\,\mathrm{d}x=\int_\Omega\tilde u(x,0)\rho(x)\,\mathrm{d}x
\end{align*}
对所有 $t\in(0,\infty)$，由 $\rho$ 的任意性得 $\tilde u(x,t)=\tilde u(x,0)$。故 $\tilde u$ 不依赖于 $t$；由等式\eqref{align:ch08-8.4} 得\eqref{align:ch08-8.8}。结论即得。
\end{proof}




\end{document}